\documentclass[11pt]{article}

\usepackage[margin=1.05in]{geometry}
\usepackage{amsmath,amssymb,amsthm}
\usepackage{mathtools}
\usepackage{graphicx}
\usepackage{booktabs}
\usepackage[colorlinks=true,allcolors=blue]{hyperref}
\usepackage[numbers,sort&compress]{natbib}

\newtheorem{theorem}{Theorem}
\newtheorem{proposition}[theorem]{Proposition}
\newtheorem{lemma}[theorem]{Lemma}
\theoremstyle{definition}
\newtheorem{definition}[theorem]{Definition}
\newtheorem{hypothesis}{Hypothesis}
\renewcommand{\thehypothesis}{H\arabic{hypothesis}}
\theoremstyle{remark}
\newtheorem{remark}[theorem]{Remark}

\newcommand{\Mstar}{M_{*}}
\newcommand{\Gcal}{\mathcal{G}}
\newcommand{\dnu}{\mathrm{d}\nu}
\newcommand{\dsig}{\mathrm{d}\sigma}
\newcommand{\Mpl}{M_{\mathrm{Pl}}}
\DeclareMathOperator{\supp}{supp}
\DeclareMathOperator{\Var}{Var}

\title{\bf Spectral structure of approximate one-form symmetry breaking\\
\large Threshold bounds from the radial breaking profile}

\author{}
\date{}

\begin{document}
\maketitle

\begin{abstract}
\noindent
The electric one-form symmetry of Maxwell theory is broken by charged matter,
and the breaking is quantified by the radial profile $\delta(r)$ of the
effective one-form charge of a Wilson line~\cite{cordova_ohmori_rudelius_2022,
basile_golmohammadi_2025}. We observe that, \emph{conditional on a positive
Stieltjes representation of the gauge-invariant static response}
(Hypothesis~\ref{hyp:stieltjes} below), the reduced profile
$\Phi(r)=\delta(r)/r^{2}$ is exactly the Laplace transform of a positive
measure $\nu$ whose support edge is the lowest threshold of the charged sector.
The word ``exact'' here means exact within that hypothesis, and nothing more:
we do \emph{not} prove that the hypothesis holds nonperturbatively, or in full
QED, where massless multi-photon cuts obstruct it. Given the representation,
standard inverse-spectral technology applies: the logarithmic slope
$\Gamma(r)=-\partial_{r}\log\Phi$ is a monotonically decreasing upper bound on
the threshold, and a local Hankel pencil refines it. We emphasise that this
machinery is not new --- it is the effective-mass and generalized-eigenvalue
technology of lattice field theory~\cite{luscher_wolff_1990,
blossier_et_al_gevp_2009}, and threshold extraction from vacuum-polarization
moments already appears in~\cite{masjuan_peris_2009}. The contribution is the
transport of that technology to a generalized-symmetry observable, together
with explicit failure modes: a no-go showing that no uniform lower bound on the
threshold exists from a finite window without a minimum-weight assumption, and
the scale invariance that obstructs Weak-Gravity inference from positivity
alone.
\end{abstract}

\section{Introduction}

Approximate global symmetries are expected to be broken in quantum
gravity~\cite{banks_seiberg_2010,harlow_ooguri_2018}, and the higher-form case
is sharp enough to be made quantitative. For the electric one-form symmetry of
Maxwell theory~\cite{gaiotto_kapustin_seiberg_willett_2015}, charged matter
screens a Wilson line, so the charge measured by a linking surface acquires a
radius dependence. C\'ordova, Ohmori and Rudelius introduced the resulting
breaking parameter and argued that its being of order one below the Planck
scale motivates the Weak Gravity
Conjecture~\cite{cordova_ohmori_rudelius_2022,arkanihamed_motl_nicolis_vafa_2006}.
Basile and Golmohammadi computed the profile explicitly at one loop for spinor
and scalar QED and for Calabi--Yau
compactifications~\cite{basile_golmohammadi_2025}.

This paper asks a narrow question about that observable: what does its
\emph{shape} determine? The answer is that, under an explicit spectral
hypothesis, the shape is a Laplace transform of a positive measure, and the
support edge of that measure is the lowest threshold that the observable
actually couples to. The inverse problem is then the familiar one of extracting
a ground-state energy from a positive Euclidean correlator.

\paragraph{What is and is not claimed.}
We are deliberate about this, because the surrounding mathematics is classical
and the adversarial literature audit accompanying this work found direct
precedents on every methodological axis.
\begin{itemize}
\item The Laplace structure of the Uehling screening correction is
      classical~\cite{uehling_1935}.
\item Complete monotonicity of a Laplace transform is Bernstein--Widder.
      Moreover, Bachas~\cite{bachas_1986} already derives $V'\ge 0$,
      $V''\le 0$ for the static potential from reflection positivity
      \emph{alone}. Our Hypothesis~\ref{hyp:stieltjes} is strictly stronger
      than reflection positivity, so the infinite hierarchy we obtain is not an
      improvement on Bachas in hypothesis-free content.
\item That $\Gamma(r)$ is a monotone upper bound converging to the support
      edge, and that a generalized eigenvalue hierarchy refines it, is the
      effective-mass/GEVP technology of lattice field
      theory~\cite{luscher_wolff_1990,blossier_et_al_gevp_2009}. We import it;
      we do not claim it.
\item Extracting threshold information from the low-energy moments of the
      vacuum polarization, by Pad\'e approximants equivalent to the same
      Hankel machinery, is done in~\cite{masjuan_peris_2009,
      bakergravesmorris1996pade}.
\end{itemize}
What remains is the assembly: the identification of the one-form breaking
profile as a completely monotone function under a stated hypothesis, the
resulting operational bounds on the charged threshold, and a careful account of
where the construction fails.

\section{The observable and the static response}
\label{sec:observable}

Work in four dimensions with $\hbar=c=1$, signature $(+,-,-,-)$, and
\begin{equation}
S=\int d^{4}x\left[-\frac{1}{4g_{R}^{2}}F_{\mu\nu}F^{\mu\nu}
+A_{\mu}j^{\mu}+\mathcal{L}_{\mathrm{mat}}\right],
\end{equation}
so that $g_R$ is fixed by the residue of the long-distance Coulomb pole. The
electric one-form current is $J_{e}=\star F/g_{R}^{2}$ and the symmetry surface
operator is $U_{\alpha}(S^{2}_{r})=\exp(i\alpha\oint_{S^{2}_{r}}\star
F/g_{R}^{2})$. Following~\cite{cordova_ohmori_rudelius_2022,
basile_golmohammadi_2025}, the effective one-form charge of a Wilson line $V$ is
\begin{equation}
q(r)=\lim_{\alpha\to0}\frac{1}{i\alpha}
\log\frac{\langle V^{\dagger}(\infty)U_{\alpha}(S^{2}_{r})V(0)\rangle}
{\langle V^{\dagger}(\infty)V(0)\rangle}.
\label{eq:qdef}
\end{equation}

\paragraph{Static potential, not a ``gauge-invariant propagator''.}
The object we constrain is the static potential between heavy probes,
\begin{equation}
V(r)=-\lim_{T\to\infty}\frac{1}{T}\log\langle W(C_{r\times T})\rangle,
\label{eq:wilson}
\end{equation}
up to the probe self-energy (perimeter) subtraction. That subtraction is
$r$-independent, so it drops from $dV/dr$ and \emph{a fortiori} from every
quantity below; no scheme ambiguity enters. We deliberately avoid any
appeal to a ``gauge-invariant gauge-field propagator'': the spectral positivity
of a gauge propagator in a covariant gauge is not a gauge-invariant statement.
Let $\Gcal(Q^{2})$ denote the three-dimensional Fourier transform of the
linear-response kernel that produces~\eqref{eq:wilson} at leading order in the
probe strength.

\begin{definition}[Breaking profile]
With $q_{\infty}=\lim_{r\to\infty}q(r)$,
\begin{equation}
\delta(r)=-\frac{r}{q_{\infty}}\,\frac{dq(r)}{dr},
\qquad
\Phi(r)=\frac{\delta(r)}{r^{2}}.
\label{eq:delta}
\end{equation}
\end{definition}
The explicit minus sign makes $\delta\ge0$ for screening: $q(r)$ decreases
towards $q_{\infty}$. Reference~\cite{basile_golmohammadi_2025} writes
$\delta_{\Lambda}=(r/q_{\infty})\,dq/dr$ without it, while their explicit
result is positive; the discrepancy is the orientation of the linking surface
and carries no physical content. We fix the convention
by~\eqref{eq:delta} and verify agreement with their explicit expression
numerically.

\section{Hypotheses}

\begin{hypothesis}\label{hyp:qft}
The response arises in a unitary, Poincar\'e-invariant quantum field theory in
$3+1$ dimensions with a stable vacuum, and rotational symmetry for the point
source. Unitarity alone does \emph{not} establish
Hypothesis~\ref{hyp:stieltjes}.
\end{hypothesis}

\begin{hypothesis}\label{hyp:gap}
The static kernel has a Coulomb pole of residue $g_{R}^{2}>0$, and its
additional spectral part is supported in $[s_{*},\infty)$ with $s_{*}>0$. A
declared lower bound on the support need not be its effective infimum.
\end{hypothesis}

\begin{hypothesis}[Positive Stieltjes representation]\label{hyp:stieltjes}
The linear-response kernel admits
\begin{equation}
\Gcal(Q^{2})=\frac{g_{R}^{2}}{Q^{2}}
+\int_{s_{*}}^{\infty}\frac{\dsig(s)}{Q^{2}+s},
\qquad \dsig\ge0,\qquad Q^{2}=|\mathbf{p}|^{2}>0,
\label{eq:H3}
\end{equation}
with $\int\dsig(s)/(\mu_{0}^{2}+s)<\infty$ for some reference scale
$\mu_{0}>0$, which is what makes the unsubtracted representation well defined.
\end{hypothesis}

\begin{hypothesis}\label{hyp:growth}
$\int s^{k}e^{-r\sqrt{s}}\dsig(s)<\infty$ for all $r>0$ and integers $k\ge0$.
Polynomial growth suffices.
\end{hypothesis}

\begin{hypothesis}\label{hyp:probe}
The source is non-dynamical and infinitely heavy, and the linear-response
relation $\tilde\varphi(\mathbf{p})=q_{W}\Gcal(Q^{2})$ holds in the regime used.
\end{hypothesis}

Hypothesis~\ref{hyp:stieltjes} is the only one with dynamical content. At order
$g_{R}^{2}$ it is not an assumption but a consequence of Lehmann positivity of
the \emph{matter} current correlator (Appendix~A). Beyond that order it is an
assumption about the Coulomb-phase screening function, and
Section~\ref{sec:scope} delimits where it fails. Hypothesis~\ref{hyp:probe}
excludes precisely the strong-source corrections of
Wichmann and Kroll~\cite{wichmann_kroll_1956}.

\section{Exact conditional Laplace representation}

\begin{theorem}[Laplace representation]\label{thm:A}
Under Hypotheses~\ref{hyp:qft}--\ref{hyp:probe}, for all $r>0$,
\begin{equation}
\Phi(r)=\int_{\Mstar}^{\infty}e^{-rx}\,\dnu(x),
\qquad \Mstar=\sqrt{s_{*}},
\label{eq:master}
\end{equation}
where $\nu$ is defined \emph{without assuming absolute continuity} as the
weighted pushforward of $\sigma$ under $x=\sqrt{s}$:
\begin{equation}
\int f(x)\,\dnu(x)\;\equiv\;\frac{1}{g_{R}^{2}}\int s\,f(\sqrt{s})\,\dsig(s).
\label{eq:pushforward}
\end{equation}
\end{theorem}

\begin{proof}
The Yukawa transform gives
$\varphi(r)=\frac{q_{W}}{4\pi r}[g_{R}^{2}+\int e^{-r\sqrt{s}}\dsig(s)]$, and
Gauss's law $q(r)=4\pi r^{2}E_{r}/g_{R}^{2}$ yields
\begin{equation}
\frac{q(r)}{q_{W}}=1+\frac{1}{g_{R}^{2}}
\int\bigl(1+r\sqrt{s}\bigr)e^{-r\sqrt{s}}\,\dsig(s).
\end{equation}
All factors of $4\pi$ and all powers of $r$ cancel. The kernel obeys the exact
identity
\begin{equation}
\frac{d}{dr}\Bigl[(1+r\sqrt{s})e^{-r\sqrt{s}}\Bigr]=-\,s\,r\,e^{-r\sqrt{s}},
\label{eq:kernelid}
\end{equation}
in which the two terms of order $\sqrt{s}$ cancel; this is what removes the
$1/s$ inherited from the subtracted dispersion relation. Differentiation under
the integral is licensed by Hypothesis~\ref{hyp:growth}, and
Hypothesis~\ref{hyp:gap} gives $q_{\infty}=q_{W}$. Inserting
into~\eqref{eq:delta} and applying~\eqref{eq:pushforward} gives
\eqref{eq:master}. An atom $w\,\delta_{s_{a}}$ in $\sigma$ maps to an atom of
weight $s_{a}w/g_{R}^{2}$ at $x=\sqrt{s_{a}}$; applying a density formula to an
atom is not permitted, and~\eqref{eq:pushforward} makes this automatic.
\end{proof}

\begin{remark}
``Exact'' means exact under Hypotheses~\ref{hyp:qft}--\ref{hyp:probe} and in
the linear-response coefficient. Higher powers of $g_{R}$ are then not an error
term but higher-order contributions to the exact measure $\sigma$. This is a
change in logical status, not a claim that Hypothesis~\ref{hyp:stieltjes} has
been proved nonperturbatively.
\end{remark}

At leading order $\dsig=g_{R}^{4}\rho_{J}(s)\,ds/s+O(g_{R}^{6})$ with $\rho_J$
the matter current spectral density, so
$\Phi(r)=g_{R}^{2}\int\rho_{J}(s)e^{-r\sqrt{s}}ds+O(g_{R}^{4})$. For one Dirac
species this reproduces~\cite[eq.~(13)]{basile_golmohammadi_2025} identically,
and the potential carries the Uehling prefactor $2\alpha/3\pi$~\cite{uehling_1935}.
The sign chain is derived from the quadratic effective action in Appendix~A and
is enforced by a test that fails if the convention is changed inconsistently.

\section{Threshold bounds}

\begin{theorem}[Logarithmic slope]\label{thm:C}
Let $\Gamma(r)=-\partial_{r}\log\Phi(r)$ and let $\langle\cdot\rangle_{r}$
denote expectation under $e^{-rx}\dnu(x)/\Phi(r)$. Then
\begin{equation}
\Gamma(r)=\langle x\rangle_{r}\ \ge\ \Mstar,
\qquad
\Gamma'(r)=-\Var_{r}(x)\ \le\ 0,
\end{equation}
and if $\Mstar=\inf\supp\nu$ then $\Gamma(r)\downarrow\Mstar$.
\end{theorem}

This is the effective mass of a positive Euclidean
correlator~\cite{luscher_wolff_1990}; the proof is the standard one and is
given in Appendix~B for completeness. Defining the local moments
$a_{n}(r)=(-1)^{n}\Phi^{(n)}(r)=\int x^{n}e^{-rx}\dnu$ and the pencil
$(H_{0})_{ij}=a_{i+j}$, $(H_{1})_{ij}=a_{i+j+1}$ for $0\le i,j\le K$:

\begin{theorem}[Local hierarchy]\label{thm:E}
When $H_{0}\succ0$,
\begin{equation}
B_{K}(r)=\lambda_{\min}(H_{1},H_{0})
=\inf_{0\neq p,\ \deg p\le K}
\frac{\int x\,p(x)^{2}e^{-rx}\dnu}{\int p(x)^{2}e^{-rx}\dnu}
\ \ge\ \Mstar,
\end{equation}
with $B_{K+1}\le B_{K}$ and $B_{0}=\Gamma$. If in addition
$\Mstar=\inf\supp\nu$, then $B_{K}\downarrow\Mstar$.
\end{theorem}

This is the GEVP hierarchy of~\cite{blossier_et_al_gevp_2009} in a radial
variable, and, through the Pad\'e--Stieltjes correspondence, the same machinery
used in~\cite{masjuan_peris_2009}. The convergence proof (Appendix~B) is
self-contained: determinacy follows from Carleman's criterion, and the approach
to the support edge from density of polynomials in $L^{2}$ of a measure
possessing an exponential moment. Carleman is used only for determinacy and
is not by itself a proof of edge convergence.

\paragraph{Edge law.}
If $\dnu/dx=C\,t^{p-1}[1+\beta t+O(t^{2})]$ with $t=x-\Mstar$, Watson's lemma
applied separately to $\Phi$ and to $\int t\,e^{-rx}\dnu$ gives
$\Gamma(r)=\Mstar+p/r+\beta p/r^{2}+O(r^{-3})$. For one Dirac species
$p=3/2$ and $\Gamma=2m+3/(2r)-5/(16mr^{2})+45/(32m^{2}r^{3})+\cdots$; for a
complex scalar $p=5/2$. The leading coefficient therefore discriminates the
channel, and both values are confirmed numerically (Section~\ref{sec:numerics}).

\section{Scope and failure modes}
\label{sec:scope}

\paragraph{Complete monotonicity is not equivalent to
Hypothesis~\ref{hyp:stieltjes}.}
By Bernstein--Widder~\cite{widder1941laplace}, $e^{\Mstar r}\Phi(r)$ is
completely monotone if and only if a positive measure with that support edge
exists. That equivalence concerns the Laplace representation \emph{only}.
Reconstructing the Stieltjes kernel~\eqref{eq:H3} additionally requires
$\int\dnu(x)/[x^{2}(\mu_{0}^{2}+x^{2})]<\infty$: taking $\dnu=x^{3}
\mathbf{1}_{x\ge1}dx$ gives a completely monotone $\Phi$ with a gap whose
integral diverges logarithmically. And $\Phi(r)=1/r$ is completely monotone
with no gap at all. Screening alone is weaker still:
$\Phi=e^{-r}-\tfrac12e^{-2r}$ is positive with positive $-\Phi'$, yet
$\Phi''<0$ for $r<\log 2$, so a screening-only diagnostic is fooled by a signed
measure. Our numerical pipeline exploits this: a positivity gate on
$(-1)^{n}\Phi^{(n)}$ and on $H_{0}$ fires before any bound is reported.

\paragraph{The first threshold need not be charged.}
In full QED the photon vacuum polarization receives contributions from
massless multi-photon intermediate states whose cut starts at
$s=0$~\cite{beneke_ruizfemenia_2016}. Hypothesis~\ref{hyp:gap} with
$s_{*}>0$ therefore cannot be inferred from the electron mass.
Theorem~\ref{thm:C} returns the infimum of the \emph{effectively coupled}
support; if that is zero, the bound is zero. A gapless model confirms this
numerically.

\begin{theorem}[No uniform lower bound]\label{thm:H}
Fix $M>\mu\ge0$, a window $0\le t\le T$ and $0<\epsilon<1$. The normalised
profiles $f_{0}=e^{-Mt}$ and
$f_{\epsilon}=(1-\epsilon)e^{-Mt}+\epsilon e^{-\mu t}$ are Laplace transforms of
positive measures with edges $M$ and $\mu$, yet
$\sup_{[0,T]}|f_{\epsilon}-f_{0}|\le\epsilon$. Hence no lower bound on $\Mstar$
holds uniformly over positive measures given finite-precision data on a finite
window.
\end{theorem}

Upper bounds survive; two-sided localisation of the edge does not, absent a
minimum-weight assumption, a spectral model, or external input. The crossover
radius is $t_{\times}=\log[(1-\epsilon)/\epsilon]/(M-\mu)$.

\paragraph{Other excluded regimes.}
In spatial dimension $d\neq3$ the kernel involves $K_{d/2-1}(xr)$ rather than
the elementary $K_{1/2}$, so the simple Laplace form does not transport by
replacing $r^{2}$. Higgs phases, confinement and antiscreening kernels require
separate treatment, and a globally antiscreening static response cannot satisfy
Hypothesis~\ref{hyp:stieltjes}.

\section{Numerical validation}
\label{sec:numerics}

All numbers are regenerated by the scripts in \texttt{reproducibility/} at
30--200 decimal digits. These are high-precision numerical checks, not interval
quadrature certificates; the distinction is maintained in the code.

Exactly solvable models are recovered exactly: a single atom gives
$\Gamma=M$ and $B_{0}=M$; two atoms are resolved at $K=1$. For one Dirac
species $r(\Gamma-2m)\to1.4981$ against the predicted $3/2$, and for a complex
scalar $\to 2.4904$ against $5/2$. The gapless model gives
$\Gamma(10^{3})=10^{-3}\to0$.

Two adversarial tests were specified in advance.

\paragraph{Signed measure.}
For $\dnu=\delta_{1}-\tfrac12\delta_{2}$ the gate fires at
$a_{3}(1)=-0.4715<0$ and $\det H_{0}=-0.1839<0$, and no bound is reported,
while $\Phi>0$ and $-\Phi'>0$ both hold --- confirming that a screening-only
test is insufficient.

\paragraph{Vanishing weight at the true edge.}
For $\dnu=10^{-12}\delta_{2}+$ continuum from $x=3$, the predicted crossover is
$r_{\times}=22.76$. Observed: $\Gamma=3.136$ at $r=10$, $3.018$ at $r=20$,
$2.532$ at $r=22.76$, $2.0005$ at $r=30$. At $r=3$ the hierarchy
$B_{0..4}$ descends only from $3.375$ to $3.108$ and never approaches $2$.
Theorem~\ref{thm:C} is never violated; it is the informativeness, not the
validity, of the bound that degrades. This is the operational content of
Theorem~\ref{thm:H} and is reported as a limitation of the method.

Conditioning is quantified rather than asserted: in a model with atoms at $2$
and $200$ the pencil is numerically singular at 60 digits and requires roughly
86 to resolve, against the 16 of double precision. High-order Hankel bounds
must not be reported from double-precision arithmetic.

\section{Obstruction to Weak-Gravity inference}

The estimator and every hierarchy bound are invariant under
$\nu\mapsto Z\nu$ for $Z>0$. The normalised geometry therefore cannot fix
$g|q|\Mpl/m$: shape determines the mass scale, while the magnitude carries
$g^{2}q^{2}$ and the density of charged states. Positivity alone cannot yield
the Weak Gravity Conjecture; a non-homogeneous gravitational input is required.

With one supplied explicitly --- a species bound $N\Lambda^{2}\le\kappa\Mpl^{2}$
\cite{dvali_species_2007}, a quantified strong-breaking threshold
$\delta_{\mathrm{obs}}(1/\Lambda)\ge\eta$, a heavy-tail bound $\tau$ and a
model-error bound $\varepsilon$, all fixed in advance --- the one-species cap
$0\le\delta_{D}\le g_{R}^{2}q^{2}/6\pi^{2}$ gives
\begin{equation}
\max_{i}\frac{g_{R}|q_{i}|\Mpl}{m_{i}}
\ \ge\ \sqrt{\frac{6\pi^{2}(\eta-\tau-\varepsilon)}{\kappa}}.
\end{equation}
This is a proved conditional implication, not a derivation of the conjecture.
Its content lies entirely in the hypotheses: an order-one breaking profile may
compromise the perturbative expansion that produced the cap, $\varepsilon$
cannot be discarded for convenience, and calibrating $\eta$ or $\kappa$ after
seeing the spectrum would make the conclusion circular.

\section{Discussion}

The analytic ingredients used here are classical, and the inverse-spectral
method is standard lattice technology. What the paper adds is that a
generalized-symmetry observable, introduced for an entirely different purpose,
carries a positive Laplace structure under a stated hypothesis, and that the
structure is strong enough to bound the charged threshold and weak enough to be
falsifiable: a measured profile violating complete monotonicity refutes
Hypothesis~\ref{hyp:stieltjes} for that observable.

The open questions are physical rather than mathematical: establishing
Hypothesis~\ref{hyp:stieltjes} for an interacting kernel beyond leading order;
constructing a positive label separating charged channels from massless cuts;
and controlling spectral tails in gravitational examples. We note finally that
Bachas' concavity bound~\cite{bachas_1986} obtains two derivative conditions on
$V(r)$ from reflection positivity with no spectral hypothesis at all. Whether
reflection positivity alone constrains higher derivatives of the radial
breaking profile --- and thus whether Hypothesis~\ref{hyp:stieltjes} can be
weakened --- we leave open.

\paragraph{Reproducibility.}
Code, data, the adversarial literature audit and a per-claim novelty matrix
accompany this work.

\appendix

\section{The vacuum-polarization sign, from the quadratic effective action}
\label{app:sign}

The sign is derived once and propagated; it is not adjusted to taste. A unit
test (\nolinkurl{tests/test_sign_convention.py}) fails if any step below is
changed inconsistently, and in particular asserts that the opposite convention
produces a \emph{negative} spectral measure and antiscreening.

Integrating out the matter sector gives, in Euclidean signature and on the
transverse subspace,
\begin{equation}
S^{(2)}=\frac12\int_{q}A_{\mu}(-q)
\Bigl[\tfrac{1}{g^{2}}\bigl(q^{2}\delta_{\mu\nu}-q_{\mu}q_{\nu}\bigr)
+\Pi_{\mu\nu}(q)\Bigr]A_{\nu}(q),
\qquad
\Pi_{\mu\nu}=\bigl(q^{2}\delta_{\mu\nu}-q_{\mu}q_{\nu}\bigr)\Pi_{E}(q^{2}).
\end{equation}
The transverse kernel is $K_{T}(Q^{2})=Q^{2}\bigl(g^{-2}+\Pi_{E}(Q^{2})\bigr)$
and the static response is its inverse,
\begin{equation}
\Gcal(Q^{2})=\frac{g^{2}}{Q^{2}\bigl[1+g^{2}\Pi_{E}(Q^{2})\bigr]}.
\label{eq:Kinv}
\end{equation}
For the ungauged matter current, fix
\begin{equation}
i\!\int\! d^{4}x\,e^{ipx}\langle T j_{\mu}(x)j_{\nu}(0)\rangle
=(p_{\mu}p_{\nu}-p^{2}\eta_{\mu\nu})\Pi(p^{2}),
\qquad
\Pi(p^{2})-\Pi(0)=p^{2}\!\int\!\frac{\rho_{J}(s)\,ds}{s(s-p^{2}-i0)},
\end{equation}
with $\rho_{J}\ge0$ by Lehmann positivity, assuming one subtraction suffices.
Continuing to spacelike $p^{2}=-Q^{2}$ and defining
\begin{equation}
\overline{\Pi}(Q^{2})\equiv-\bigl[\Pi(-Q^{2})-\Pi(0)\bigr]
=Q^{2}\int\frac{\rho_{J}(s)\,ds}{s(s+Q^{2})}\ \ge\ 0,
\end{equation}
which is non-negative and increasing, and renormalising $\Pi_{E}(0)=0$ so that
$g=g_{R}$ is the long-distance coupling, \eqref{eq:Kinv} becomes
\begin{equation}
\Gcal(Q^{2})=\frac{g_{R}^{2}}{Q^{2}\bigl[1-g_{R}^{2}\overline{\Pi}(Q^{2})\bigr]}
=\frac{g_{R}^{2}}{Q^{2}}
+g_{R}^{4}\int\frac{\rho_{J}(s)\,ds}{s(s+Q^{2})}+O(g_{R}^{6}),
\label{eq:expansion}
\end{equation}
so that
\begin{equation}
\dsig_{\mathrm{LO}}(s)=g_{R}^{4}\,\frac{\rho_{J}(s)}{s}\,ds\ \ge\ 0 ,
\end{equation}
which is Hypothesis~\ref{hyp:stieltjes} at leading order. The next order does
not inherit this. The $O(g_R^6)$ term of~\eqref{eq:expansion} is
$g_R^6Q^2\bigl[\int\rho_J\,ds/(s(s+Q^2))\bigr]^2$, which vanishes at $Q^2=0$,
rises and then falls like $1/Q^2$. A function $\int\dsig/(Q^2+s)$ with
$\dsig\ge0$ is largest at $Q^2=0$, so this term is not of that form, and
truncating~\eqref{eq:expansion} at fixed order does not preserve the
hypothesis. Whatever positivity the chain has belongs to its sum.

\paragraph{The resummed bubble chain.}
Writing $Q^2/[s(s+Q^2)]=1/s-1/(s+Q^2)$ puts the denominator
of~\eqref{eq:expansion} in the form
\begin{equation}
W(Q^2)\equiv1-g_R^2\overline{\Pi}(Q^2)
=Z_3+g_R^2\int\frac{\rho_J(s)\,ds}{s+Q^2},
\qquad
Z_3\equiv1-g_R^2\int\frac{\rho_J(s)}{s}\,ds ,
\label{eq:Wpf}
\end{equation}
so that $W(0)=1$, $W(Q^2)\to Z_3$ as $Q^2\to\infty$, and $W$ decreases
strictly on $Q^2>0$ whenever $\rho_J\ge0$ does not vanish identically.
The sign of $Z_3$ decides the outcome.

If $Z_3>0$, $W/g_R^2$ is a Stieltjes function. By the duality between Stieltjes
and complete Bernstein functions~\cite{schilling_song_vondracek2012},
$Q^2W/g_R^2=1/\Gcal$ is then complete Bernstein and $\Gcal$ is Stieltjes. Its
Coulomb residue is $g_R^2/W(0)=g_R^2$, and $0<\Gcal\le g_R^2/(Q^2Z_3)$ excludes
an additive constant, so the resummed chain satisfies
Hypothesis~\ref{hyp:stieltjes} to all orders in the chain.

If $Z_3<0$, $W$ has exactly one zero, at some $Q^2_{\mathrm g}>0$. There $\Gcal$
has a simple pole with residue $g_R^2/[Q^2_{\mathrm g}W'(Q^2_{\mathrm g})]<0$, at
spacelike momentum, which no Stieltjes function admits.

If $Z_3=0$ there is no such pole, but
$\Gcal(Q^2)=\bigl[\int Q^2\rho_J\,ds/(s+Q^2)\bigr]^{-1}\to1/\int\rho_J\,ds$,
a positive constant whenever that total mass is finite, and the hypothesis
excludes a constant. The absence of a spacelike pole is therefore necessary for
Hypothesis~\ref{hyp:stieltjes} and not sufficient.

Two consequences bear on the rest of the paper. The continuum spectral density
of the chain is $g_R^4\rho_J(s)/(s\,|W(-s-i0)|^2)$, non-negative for every
coupling, $Z_3<0$ included. A positivity test built from moments of the
continuum therefore cannot see the failure at $Z_3<0$, which lives in a discrete
pole; the finite-moment conditions of Section~\ref{sec:scope} have exactly this
limitation. Second, if $\rho_J$ is cut off sharply at $s=\Lambda^2$ and does not
vanish there, as the Dirac density does not, then $W$ is real again for
$Q^2<-\Lambda^2$ and, when $Z_3>0$, vanishes once in that range. The resummed
measure then carries an atom of positive weight above the cutoff, outside the
support of $\rho_J$. This is the mechanism studied
in~\cite{giacosa_wolkanowski_2012}, where resummation places a pole on the
physical sheet outside the input spectrum, with positive residue, and a
spectral sum rule restores the normalisation. Theorem~\ref{thm:A} already
admits atoms in $\sigma$ through~\eqref{eq:pushforward}.

For the one-loop Dirac density at $\alpha=1/137$, $Z_3=0.966$ at
$\Lambda^2=10^{20}\,m^2$, and $Z_3$ reaches zero only near
$\log(\Lambda^2/4m^2)=3\pi/\alpha$, the Landau pole of the one-loop coupling.
Below it the chain satisfies the hypothesis with a wide margin. None of this
bears on the irreducible $O(g_R^6)$ contributions or on the nonperturbative
kernel, where Hypothesis~\ref{hyp:stieltjes} remains an assumption.

\paragraph{Check 1 (inverse kernel).}
The $O(g_{R}^{4})$ coefficient in \eqref{eq:expansion} is
$+\rho_{J}/[s(s+Q^{2})]$. Writing $1+g_{R}^{2}\overline{\Pi}$ in place of
$1-g_{R}^{2}\overline{\Pi}$ --- with $\overline{\Pi}\ge0$ --- flips it to
$-\rho_{J}/[s(s+Q^{2})]$ and destroys the hypothesis. This is the
discriminating check.

\paragraph{Check 2 (Uehling).}
Inserting the one-loop Dirac density
$\rho_{J}=\frac{q^{2}}{12\pi^{2}}(1+2m^{2}/s)\sqrt{1-4m^{2}/s}$ and
substituting $s=4m^{2}\zeta^{2}$ gives the potential correction prefactor
$g_{R}^{2}q^{2}/6\pi^{2}$, which for $q=1$ and $g_{R}^{2}=4\pi\alpha$ equals
$2\alpha/3\pi$: the Uehling coefficient~\cite{uehling_1935}.

\paragraph{Check 3 (screening).}
$g^{2}_{\mathrm{eff}}(Q^{2})=Q^{2}\Gcal(Q^{2})=g_{R}^{2}/[1-g_{R}^{2}
\overline{\Pi}]$ is increasing in $Q^{2}$, so the coupling grows in the
ultraviolet. The flipped convention gives a decreasing effective coupling,
i.e. antiscreening --- the signature of the error.

\paragraph{One-species cap.}\label{app:cap}
For the one-loop Dirac density, put $u=4m^2/x^2\in[0,1]$ and
$f(u)=(1+u/2)\sqrt{1-u}$. Since
$f(u)^2=1-3u^2/4-u^3/4$, one has $0\le f(u)\le1$. Consequently
\begin{equation}
0\le\delta_D(r)=\frac{g_R^2q^2}{6\pi^2}r^2
\int_{2m}^{\infty}x e^{-rx}f(4m^2/x^2)\,dx
\le\frac{g_R^2q^2}{6\pi^2}r^2\int_0^{\infty}x e^{-rx}\,dx
=\frac{g_R^2q^2}{6\pi^2}.
\end{equation}
This bounds the one-loop coefficient, not the uncomputed interacting remainder.

\section{Proofs}
\label{app:proofs}

\paragraph{Theorem~\ref{thm:C}.}
Write $d\mu_{r}=e^{-rx}\dnu/\Phi(r)$, a probability measure. Then
$\Gamma=-\Phi'/\Phi=\langle x\rangle_{r}\ge\Mstar$ since $\nu$ is supported in
$[\Mstar,\infty)$, and differentiating the ratio of moments gives
$\Gamma'=-(\langle x^{2}\rangle_{r}-\langle x\rangle_{r}^{2})=-\Var_{r}(x)$.
For convergence, fix $\varepsilon>0$ and $r_{0}>0$. The interval
$I=[\Mstar,\Mstar+\varepsilon/4]$ has mass $c>0$ when $\Mstar=\inf\supp\nu$, so
$\Phi(r)\ge c\,e^{-r(\Mstar+\varepsilon/4)}$. In the numerator of
$\Gamma-\Mstar$ the contribution from $x\le\Mstar+\varepsilon/2$ is at most
$(\varepsilon/2)\Phi$, and the tail is bounded by
$A\,e^{-(r-r_{0})(\Mstar+\varepsilon/2)}$ with
$A=\int(x-\Mstar)e^{-r_{0}x}\dnu<\infty$ by Hypothesis~\ref{hyp:growth}. The
ratio of the tail to $\Phi$ is $O(e^{-r\varepsilon/4})\to0$, so
$\limsup_{r}(\Gamma-\Mstar)\le\varepsilon/2$; letting $\varepsilon\to0$ and
using monotonicity gives the limit. \qed

\paragraph{Theorem~\ref{thm:E}, bound and monotonicity.}
For $p_{v}(x)=\sum_{i\le K}v_{i}x^{i}$,
\begin{equation}
v^{\mathsf{T}}(H_{1}-\Mstar H_{0})v
=\int(x-\Mstar)\,p_{v}(x)^{2}e^{-rx}\dnu\ \ge\ 0,
\end{equation}
since $x\ge\Mstar$ on the support; the Rayleigh characterisation and the bound
follow. Enlarging the polynomial space from degree $K$ to $K+1$ can only lower
an infimum, giving $B_{K+1}\le B_{K}$; and $K=0$ gives $a_{1}/a_{0}=\Gamma$.
$H_{0}\succ0$ provided the support contains at least $K+1$ points; for a
measure with $L$ atoms the pencil is exact at order $L-1$ and singular
beyond, which must be handled by rank reduction rather than reported as a new
bound. \qed

\paragraph{Theorem~\ref{thm:E}, convergence.}
Fix $r>0$. First, \emph{determinacy}: for $0<t<r$, using
$e^{tx}\ge(tx)^{2n}/(2n)!$,
\begin{equation}
a_{2n}(r)\le(2n)!\,t^{-2n}\,\Phi(r-t),
\qquad\text{hence}\qquad
\sum_{n\ge1}a_{2n}(r)^{-1/(2n)}=\infty,
\end{equation}
which is Carleman's criterion~\cite{akhiezer1965moment}. Carleman establishes
determinacy and nothing more; it is not by itself a proof of edge convergence.

Second, \emph{edge convergence}. Let $d\omega=(1+x)e^{-rx}\dnu$, which is
finite and possesses an exponential moment, $\int e^{tx}d\omega<\infty$ for
$0<t<r$. Polynomials are therefore dense in $L^{2}(\omega)$: if $f\perp$ all
polynomials, Cauchy--Schwarz makes $z\mapsto\int f(x)e^{zx}d\omega$ analytic in
a strip about the imaginary axis with all derivatives vanishing at the origin,
so it vanishes identically and $f=0$ by uniqueness of the Fourier transform.
Now let $\varepsilon>0$ and approximate in $L^{2}(\omega)$ the indicator of
$[\Mstar,\Mstar+\varepsilon]$, which has positive mass when
$\Mstar\in\supp\nu$. Because $d\omega$ dominates both $e^{-rx}\dnu$ and
$x\,e^{-rx}\dnu$, the $L^{2}(\omega)$ norm controls both integrals in the
Rayleigh quotient, so the quotient of the approximating polynomials converges
to that of the indicator, which is at most $\Mstar+\varepsilon$. Hence
$\lim_{K}B_{K}\le\Mstar+\varepsilon$ for every $\varepsilon>0$; with
$B_{K}\ge\Mstar$ the claim follows. Finally $\supp\mu_{r}=\supp\nu$ because
$e^{-rx}>0$ everywhere, so the edge reached is $\inf\supp\nu$. \qed

\paragraph{Edge law.}
Apply Watson's lemma~\cite{dlmf} separately to $\Phi(r)=\int e^{-rx}\dnu$ and
to $\int t\,e^{-rx}\dnu$ with $t=x-\Mstar$, then form
$\Gamma-\Mstar=(\int t\,e^{-rx}\dnu)/\Phi$. Subtracting $\Mstar\Phi$ before
expanding avoids differentiating an asymptotic series. With
$\dnu/dx=Ct^{p-1}[1+\beta t+\gamma t^{2}+O(t^{3})]$ and
$\int_{0}^{\infty}t^{p+j-1}e^{-rt}dt=\Gamma_{E}(p+j)r^{-p-j}$, dividing the two
series gives
\begin{equation}
\Gamma(r)=\Mstar+\frac{p}{r}+\frac{\beta p}{r^{2}}
+\frac{2\gamma p(p+1)-\beta^{2}p^{2}}{r^{3}}+O(r^{-4}).
\end{equation}
If only $1+O(t)$ is known, only $\Gamma=\Mstar+p/r+O(r^{-2})$ follows. For an
atom $Z\delta_{\Mstar}$ \emph{separated} from the rest of the support by
$d>0$, convergence is $O(e^{-dr})$; separation is a hypothesis, not a
consequence of the atom. If an atom and a continuum share the same edge,
$\Gamma-\Mstar\sim C\Gamma_{E}(p+1)r^{-p-1}/Z$: algebraic despite the pole.
\qed


\bibliographystyle{JHEP}
\bibliography{references}

\end{document}
